\chapter{KẾT LUẬN VÀ HƯỚNG PHÁT TRIỂN}
\label{chap:chapter5}

\section{Tổng kết các kết quả đạt được}
\label{sec:tongKet}

Khóa luận đã hoàn thành bốn mục tiêu nghiên cứu chính và thực hiện thêm một
nhóm phân tích chuyên sâu (mục thứ năm dưới đây):

\begin{enumerate}
    \item Xây dựng môi trường mô phỏng: phát triển thành công
    \path{MultiWarehouseInventoryEnv} theo chuẩn Gymnasium
    \parencite{towers2024gymnasium}, mô hình hóa đầy đủ hệ thống $10$ kho,
    $30$ mặt hàng ($300$ cặp tác tử), với nhu cầu ngẫu nhiên, thời gian giao
    hàng biến động $[1, 3]$ ngày, ràng buộc sức chứa dùng chung tính riêng cho
    từng kho theo quy mô nhu cầu của kho đó, bốn thành phần chi phí vận hành
    (lưu kho, thiếu hàng, đặt hàng, tràn kho) và hình phạt mức phục vụ tách
    riêng khỏi chi phí vận hành. Nhu cầu lấy từ bộ dữ liệu M5 Walmart
    \parencite{makridakis2022background}.

    \item Thiết kế hàm phần thưởng: phần thưởng phân rã cục bộ (local
    factored reward) theo từng cặp kho--mặt hàng, với động cơ giảm khó khăn gán
    tín dụng \parencite{wolpert2001}, kèm hình phạt mức phục vụ dạng ràng buộc
    mềm tính trên cửa sổ trượt $30$ ngày.

    \item Cài đặt và huấn luyện IPPO: triển khai thuật toán Independent
    Multi-Agent PPO \parencite{dewitt2020ippo} với cơ chế chia sẻ tham số
    \parencite{gupta2017}, sử dụng GAE \parencite{schulman2016gae} và chuẩn hóa
    lợi thế theo từng cặp. Mô hình hội tụ ổn định trên quy mô $300$ tác
    tử sau khoảng $600$--$800$ episode, với entropy giảm từ $1{,}778$ xuống
    $0{,}18$ và explained variance của Critic duy trì quanh $0{,}75$--$0{,}80$,
    nhất quán trên cả $3$ hạt giống huấn luyện độc lập đã kiểm chứng.

    \item Đánh giá đối chứng: so sánh IPPO với ba chính sách cổ điển
    (EOQ, $(s, S)$, Newsvendor) trên $30$ episode đánh giá dùng chung hạt
    giống, lặp lại trên $3$ hạt giống huấn luyện độc lập, có kiểm định thống kê
    \parencite{henderson2018, agarwal2021}, và bổ sung phép so sánh ở cùng mức
    phục vụ nhằm loại bỏ sai lệch do các phương pháp nằm ở những điểm khác
    nhau trên đường đánh đổi.

    \item Phân tích theo giai đoạn và kiểm tra hành vi: tách kết quả theo giai
    đoạn nhu cầu (ổn định, biến động, sự kiện, mùa lễ) và thử sốc cầu nhân tạo
    (thực nghiệm B); kiểm tra bằng thực nghiệm tính hợp lý của hành vi đặt
    hàng, dấu hiệu khai thác phần thưởng và tầm quan trọng của từng nhóm đặc
    trưng trạng thái. Kết quả mô phỏng được đưa lên ứng dụng Streamlit cho
    phép xem diễn biến theo từng ngày của từng cặp kho--mặt hàng (nhu cầu, tồn
    kho, hàng đang về, lượng đặt, thiếu hàng, phần thưởng, cảnh báo thiếu hàng).
\end{enumerate}

\subsection{Trả lời các câu hỏi nghiên cứu}
\label{subsec:traLoiCauHoi}

Câu hỏi 1 (IPPO so với các chính sách cổ điển). Câu trả lời là
khẳng định, nhưng chỉ khi mức phục vụ được kiểm soát. Ở điểm vận hành mà mỗi
phương pháp tự chọn, IPPO \textbf{không} rẻ hơn các chính sách cổ điển: gộp
trên $3$ hạt giống huấn luyện $\times$ $30$ episode đánh giá, IPPO đắt hơn
$(s, S)$ $4{,}89\%$ và Newsvendor $3{,}16\%$ (dù vẫn rẻ hơn EOQ $10{,}41\%$),
khác biệt có ý nghĩa thống kê rất cao ($p = 5{,}45 \times 10^{-33}$ với kiểm
định $t$ Welch; Cohen's $d = 4{,}06$) và nhất quán trên cả $3$ hạt giống độc
lập. Nguyên nhân là IPPO dịch chuyển tới mức phục vụ $90{,}5\%$ --- cao hơn hẳn
$75$--$84\%$ của ba chính sách cổ điển --- nên đây không phải là ưu thế Pareto
mà là một điểm khác trên đường đánh đổi chi phí--mức phục vụ.

Khi so sánh ở cùng mức phục vụ ($\approx 90\%$) --- phép so sánh chặt chẽ hơn
về mặt phương pháp luận vì loại bỏ được sai lệch do hai mục tiêu tối ưu khác
nhau --- kết luận đảo chiều: IPPO rẻ hơn $(s, S)$ $10{,}2\%$ và rẻ hơn
Newsvendor $14{,}4\%$, trong khi EOQ không có cấu hình nào đạt được mức phục vụ
đó do cấu trúc lượng đặt hàng cố định. Nguyên nhân là các chính sách cổ điển
chỉ có thể nâng mức phục vụ bằng cách tăng tham số an toàn áp dụng đồng loạt
cho cả $300$ cặp --- khiến chi phí tăng đáng kể ở cả những mặt hàng vốn đã đạt mức
phục vụ tốt --- còn IPPO điều chỉnh theo từng trạng thái quan sát được. Khóa
luận coi phép so sánh ở cùng mức phục vụ là căn cứ chính đáng để trả lời Câu
hỏi 1, và ghi nhận rõ ràng rằng phép so sánh trực tiếp không kiểm soát mức
phục vụ có thể dẫn tới kết luận ngược lại.

Thực nghiệm B bổ sung sắc thái cho câu trả lời này. Khi nhu cầu ổn định,
$(s,S)$ tinh chỉnh trực tiếp rẻ hơn IPPO vài phần trăm; ưu thế của IPPO tập
trung ở giai đoạn biến động --- fill rate chỉ giảm $2{,}4$ điểm phần trăm từ
ngày ổn định sang ngày biến động (so với $5$--$7$ điểm của $(s,S)$), rẻ hơn
khoảng $10\%$ ở ngày có sự kiện và cuối tuần dù phục vụ tốt hơn, và tốn ít hơn
khoảng $29\%$--$32\%$ chi phí khi chống đỡ một cú sốc tăng cầu $50\%$. Các
chênh lệch ở ngày ổn định, ngày có sự kiện và cuối tuần đều có ý nghĩa thống kê
theo bootstrap khối tuần. Mùa lễ cuối năm là giai đoạn IPPO yếu nhất, nhưng dữ
liệu chỉ có $9$ tuần nên chưa đủ để kết luận chắc chắn. Ở cùng mức phục vụ, IPPO
rẻ hơn $(s,S)$ ở cả tám nhóm ngày, và rẻ hơn ở cả $36$ bộ đơn giá chi phí trong
phân tích độ nhạy. Kết luận hợp lý vì vậy không phải ``học
tăng cường tốt hơn'' một cách chung chung, mà là: chính sách cổ điển đủ tốt khi
nhu cầu ổn định, còn IPPO có giá trị khi nhu cầu biến động nhanh.

Câu hỏi 2 (hiệu năng theo nhóm quy mô nhu cầu). Câu trả lời là có điều
kiện. Một bộ trọng số duy nhất điều khiển được $300$ cặp có nhu cầu từ dưới
$0{,}1$ đến $143$ đơn vị mỗi ngày, và ở nhóm nhu cầu trung bình và cao IPPO vượt
$(s,S)$ trên cả fill rate lẫn chi phí. Tuy nhiên, ở nhóm bán chậm (dưới $2$ đơn
vị mỗi ngày) IPPO phục vụ kém hơn $(s,S)$; phân tích phân bố cho thấy $39\%$ số
cặp của IPPO dưới ngưỡng $85\%$. Nguyên nhân chủ yếu nằm ở cách định cỡ hình
phạt mức phục vụ theo $\bar{D}^{(i)}$ chứ không phải ở khả năng biểu diễn của
mạng dùng chung (Mục~\ref{subsec:tongQuatHoa}, \ref{subsec:kiemTraStateReward}).
Khả năng áp dụng cho mặt hàng chưa từng gặp khi huấn luyện chưa được kiểm chứng.

Câu hỏi 3 (thiết kế tín hiệu phần thưởng). Về khía cạnh chuẩn hóa, có bằng chứng
sơ bộ ủng hộ. Một lần chạy loại trừ tắt chuẩn hóa phần thưởng theo cặp (dừng ở
khoảng $625$ episode) cho tổng chi phí cao gấp khoảng $3{,}2$ lần và fill rate
chỉ $80{,}7\%$, so với lần huấn luyện gốc ở cùng khoảng episode
(Mục~\ref{subsec:ablationReward}). Đáng chú ý là Critic khi đó đạt explained
variance gần $1$ nhưng chính sách vẫn kém. Điều này khớp với lập luận rằng
return chênh nhau nhiều bậc độ lớn khiến tín hiệu học bị các cặp nhu cầu lớn chi
phối. Kết luận định lượng đầy đủ chờ lần chạy $1.000$ episode trong đợt thực
nghiệm bổ sung.

\section{Những hạn chế còn tồn tại}
\label{sec:hanChe}

\begin{itemize}
    \item Cấu trúc mạng lưới đơn giản: các kho hoạt động song song và
    độc lập về nguồn cung; chưa xét cấu trúc đa cấp (multi-echelon) với quan hệ
    cung ứng giữa các kho \parencite{clark1960} và hiệu ứng bullwhip
    \parencite{lee1997}.

    \item Không có cơ chế phối hợp tường minh: các tác tử ra quyết
    định hoàn toàn dựa trên quan sát cục bộ, phối hợp chỉ diễn ra gián tiếp qua
    tín hiệu sức chứa kho. Khóa luận chưa định lượng được phần giá trị mà một
    critic tập trung theo hướng CTDE \parencite{yu2022mappo} có thể mang lại ---
    đặc biệt đáng lưu ý vì Zhu và Wu \cite{zhu2026spatiotemporal} báo cáo mức
    suy giảm đáng kể khi loại bỏ cơ chế phối hợp trong khung của họ.

    \item Số episode mỗi hạt giống huấn luyện: kết quả chính được kiểm
    chứng trên $3$ hạt giống độc lập theo khuyến nghị của Henderson và cộng sự
    \cite{henderson2018}, mỗi hạt giống $1.000$ episode. Huấn luyện tiếp hạt
    giống $42$ tới khoảng $4.083$ episode đã giảm thêm khoảng $5\%$ chi phí,
    cho thấy $1.000$ episode chưa phải điểm bão hòa; kết quả dài hơn mới có trên
    một hạt giống.

    \item Mất cân bằng mức phục vụ giữa các mặt hàng: hình phạt mức phục
    vụ tỷ lệ với nhu cầu trung bình nên chính sách phân bổ mức phục vụ thấp hơn cho
    mặt hàng bán chậm mà vẫn đạt mục tiêu chung; $39\%$ số cặp có fill rate dưới $85\%$.

    \item Giới hạn của thực nghiệm theo giai đoạn: miền test chỉ dài $491$
    ngày nên các episode chồng lấn nhau. Bootstrap theo khối tuần cho khoảng tin
    cậy của từng giai đoạn nhưng không loại bỏ hết tự tương quan, và mùa lễ chỉ
    có $9$ tuần nên kết luận về mùa lễ chưa có ý nghĩa thống kê.

    \item Kết quả chính và các phân tích bổ sung dùng hai checkpoint khác
    nhau: các bảng chính dùng $3$ hạt giống $\times$ $1.000$ episode, còn thực
    nghiệm B, phần kiểm tra hành vi và phân tích độ nhạy dùng checkpoint hạt
    giống $42$ huấn luyện tiếp tới khoảng $4.083$ episode. Đợt thực nghiệm bổ
    sung (Mục~\ref{subsec:dotBoSung}) thống nhất lại toàn bộ trên một cấu hình.

    \item Hình phạt mức phục vụ tỷ lệ theo nhu cầu: mọi kết quả ở
    Chương~\ref{chap:chapter4} dùng cách định cỡ cũ. Biến thể chuẩn
    hóa~(\ref{eq:phatChuanHoa}) đã cài đặt nhưng chưa có kết quả huấn luyện.

    \item Đặc trưng lịch gần như không được khai thác: $27/43$ chiều
    quan sát chỉ đóng góp khoảng $3{,}5\%$ chi phí theo phân tích permutation
    importance, và IPPO yếu nhất ở mùa lễ.

    \item Không gian hành động rời rạc thô: chỉ có $K = 6$ mức đặt
    hàng; trong thực tế, doanh nghiệp có thể cần quyết định lượng đặt chính xác
    hơn.

    \item Chưa tích hợp dự báo nhu cầu: tác tử sử dụng trực tiếp nhu
    cầu lịch sử thay vì dự báo bằng mô hình học sâu chuỗi thời gian.

    \item Thực nghiệm trên môi trường mô phỏng: chưa được kiểm chứng
    trên hệ thống ERP thực tế hoặc dữ liệu vận hành trực tiếp.
\end{itemize}

\section{Hướng phát triển trong tương lai}
\label{sec:huongPhatTrien}

\subsection{Đợt thực nghiệm bổ sung đã cài đặt}
\label{subsec:dotBoSung}

Các hạn chế nêu ở Mục~\ref{sec:hanChe} đều đã được chuyển thành thí nghiệm cụ
thể. Toàn bộ mã nguồn đã được cài đặt và kiểm thử (tổng cộng $39$ kiểm thử đơn
vị), và được gom vào một kịch bản duy nhất (\texttt{scripts/campaign.py}) có
khả năng huấn luyện tiếp khi bị gián đoạn. Riêng phần huấn luyện (khoảng $25$
giờ CPU) chưa hoàn tất tại thời điểm viết báo cáo. Mỗi thí nghiệm loại trừ chỉ
khác cấu hình tham chiếu đúng một thay đổi (Bảng~\ref{tab:dotBoSung}); cấu hình
tham chiếu dùng hình phạt mức phục vụ chuẩn hóa~(\ref{eq:phatChuanHoa}).

\begin{table}[htbp]
    \centering
    \caption{Đợt thực nghiệm bổ sung: mã nguồn đã sẵn sàng, đang chờ huấn luyện}
    \label{tab:dotBoSung}
    \small
    \begin{tabular}{|p{3.0cm}|p{5.6cm}|p{4.8cm}|}
        \hline
        \textbf{Lần chạy} & \textbf{Thay đổi so với tham chiếu} & \textbf{Trả lời câu hỏi} \\
        \hline
        Mô hình chính ($3$ hạt giống $\times$ $4.000$ episode) & Hình phạt mức phục vụ chuẩn hóa & Mặt hàng bán chậm còn bị phân bổ mức phục vụ thấp hơn không; kết luận ``không đắt hơn $(s,S)$'' có lặp lại trên $3$ hạt giống không \\ \hline
        Tham chiếu $1.000$ episode & --- & Mốc so sánh cho các thí nghiệm loại trừ \\ \hline
        Phần thưởng toàn cục ($3$ hạt giống) & \texttt{reward\_mode: global} & Câu hỏi 3: cục bộ so với toàn cục \\ \hline
        Tắt chuẩn hóa theo cặp ($3$ hạt giống) & \texttt{normalize\_reward\_per\_pair: false} & Câu hỏi 3: có/không chuẩn hóa \\ \hline
        Baseline theo nhóm quy mô cầu & Mỗi nhóm cầu một bộ tham số (s,S), Newsvendor, EOQ & So sánh với đối thủ mạnh hơn \\ \hline
        Độ bền theo điều kiện vận hành & Sức chứa 4/6 ngày, lead time 2--3 hoặc luôn 3 ngày, nhu cầu $\pm 20\%$ (không huấn luyện lại) & Kết luận có đứng vững khi môi trường thay đổi \\ \hline
        Phạt SLA kiểu cũ & \texttt{service\_penalty\_mode: demand} & Tác động của việc định cỡ lại hình phạt \\ \hline
        Actor--Critic chung thân & \texttt{shared\_trunk: true}, một optimizer & Vì sao phải tách hai mạng (Mục~\ref{subsec:ablationTrunk}) \\ \hline
        Bỏ đặc trưng lịch / cấp kho & \texttt{obs\_drop} & Trạng thái có biến dư thừa không \\ \hline
        Thêm sự kiện sắp tới & $2$ đặc trưng: số ngày tới sự kiện kế tiếp, có sự kiện trong $7$ ngày & Khắc phục điểm yếu mùa lễ \\ \hline
        Warm-start & Lịch sử nhu cầu khởi tạo bằng dữ liệu thật & Ảnh hưởng của lịch sử rỗng đầu episode \\ \hline
        Hold-out mặt hàng & Huấn luyện $24$ mặt hàng, đánh giá $6$ mặt hàng chưa gặp & Câu hỏi 2 (tổng quát hóa) \\
        \hline
    \end{tabular}
\end{table}

Giao thức đánh giá của đợt này cũng được củng cố: baseline tinh chỉnh trên miền
val (cùng miền dùng để chọn checkpoint IPPO) thay vì miền train; báo cáo kiểm
định $t$ theo cặp cùng hệ số $d_z$; đánh giá thêm trên các cửa sổ cố định trải
đều miền test; và dùng khoảng tin cậy bootstrap theo khối tuần cho từng giai
đoạn nhu cầu. Các công cụ phân tích này đã được áp dụng ngay cho kết quả hiện
có ở Chương~\ref{chap:chapter4}.

\subsection{Các hướng mở rộng khác}

\begin{itemize}
    \item Lead time dài hơn $3$ ngày và mục tiêu mức phục vụ khác $85\%$
    (chẳng hạn $80\%$ hoặc $90\%$): phân tích độ bền hiện tại chỉ đánh giá được
    các thay đổi mà chính sách đã học vẫn nhận làm đầu vào (sức chứa, lead time
    trong khoảng $1$--$3$ ngày, mức nhu cầu). Số chiều quan sát phụ thuộc lead
    time tối đa, còn hệ số phạt được hiệu chỉnh cho mục tiêu $85\%$, nên hai
    trường hợp này đòi hỏi huấn luyện lại cho từng cấu hình.

    \item Nhân tử Lagrange riêng cho từng cặp: thay vì cố định hệ số
    $\phi_{dv}$, cập nhật một nhân tử cho mỗi cặp trong lúc huấn luyện để ràng
    buộc mức phục vụ được thỏa mãn chặt ở từng mặt hàng.

    \item Điều chuyển hàng liên kho: bổ sung hành động chuyển hàng giữa
    các kho kèm chi phí điều chuyển vào hàm phần thưởng; khi đó không gian hành
    động và bảng đặc tả tác tử cần được mở rộng tương ứng.
    \item Mở rộng sang cấu trúc đa cấp (multi-echelon): mô hình hóa
    quan hệ cung ứng giữa kho trung tâm và kho khu vực, cho phép tác tử học
    chiến lược điều phối liên kho.

    \item So sánh với kiến trúc huấn luyện tập trung --- thực thi phi
    tập trung: bổ sung một nhánh MAPPO \parencite{yu2022mappo} trong đó Critic
    nhận thêm trạng thái ở cấp kho, nhằm đo lường trực tiếp giá trị của phối
    hợp tường minh so với phối hợp gián tiếp qua tín hiệu sức chứa. Đây là
    hướng mở rộng trực tiếp nhất về mặt cài đặt, vì kiến trúc hiện tại đã tách
    riêng mạng Actor và Critic; kết quả sẽ cho phép định lượng rõ giá trị của
    phối hợp tường minh trong cấu hình nghiên cứu này.

    \item Tích hợp cơ chế truyền thông giữa tác tử: áp dụng các
    phương pháp như CommNet hoặc TarMAC để cho phép tác tử chia sẻ thông tin,
    đặc biệt về tình trạng sức chứa kho.

    \item Không gian hành động liên tục: sử dụng PPO cho hành động liên
    tục hoặc mạng hành động phân cấp (hierarchical action) để quyết định lượng
    đặt hàng chính xác.

    \item Gán công theo độ trễ giao hàng: chi phí phát sinh ở ngày
    $t$ phần lớn là hệ quả của quyết định đặt hàng ở ngày $t - L$. Hiện việc
    gán công này được phó thác cho GAE với $\gamma = 0{,}99$; một cơ chế liên
    kết phần thưởng trực tiếp với hành động sinh ra nó, theo hướng của Nam và
    cộng sự \cite{jeong2026alca}, có thể làm tín hiệu học sắc nét hơn.

    \item Tích hợp dự báo nhu cầu: nhúng mô hình dự báo chuỗi thời
    gian (LSTM, Transformer) vào vector quan sát của tác tử, tận dụng thông tin
    xu hướng và mùa vụ. Cần lưu ý rằng kết quả hiện tại cho thấy phương án
    không dùng dự báo đã đủ vượt các chính sách cổ điển với chi phí tính toán
    thấp hơn nhiều bậc, nên hướng này cần được đánh giá trên tương quan
    lợi ích --- chi phí chứ không mặc nhiên là một cải tiến.

    \item Triển khai thử nghiệm thực tế: xây dựng hệ thống dạng SaaS
    kết nối API với nền tảng quản lý kho, cho phép đánh giá mô hình trên dữ liệu
    vận hành trực tiếp.

    \item Mở rộng quy mô: thử nghiệm với toàn bộ $3.049$ SKU gốc của
    M5 và đánh giá khả năng mở rộng của cơ chế chia sẻ tham số.

    \item Lớp giao tiếp ngôn ngữ qua MCP: dùng Model Context Protocol để
    một mô hình ngôn ngữ lớn gọi engine IPPO như một công cụ, cho phép người
    quản lý kho đặt câu hỏi bằng ngôn ngữ tự nhiên (ví dụ mặt hàng nào sắp
    thiếu) và nhận đề xuất đặt hàng. IPPO vẫn là nơi ra quyết định định lượng;
    MCP chỉ lo giao tiếp và điều phối công cụ. Một prototype minh họa luồng
    điều phối này được trình bày ở Phụ lục~\ref{app:mcp}; hướng này không
    thuộc đóng góp chính của khóa luận.
\end{itemize}